\documentclass[12pt,a4paper]{article}

% ====== ЯЗЫК И ШРИФТЫ (XeLaTeX) ======
\usepackage{fontspec}
\usepackage{polyglossia}
\setdefaultlanguage{russian}
\setotherlanguage{english}

\setmainfont{CMU Serif}
\setsansfont{CMU Sans Serif}
\setmonofont{CMU Typewriter Text}
\newfontfamily\cyrillicfont{CMU Serif}
\newfontfamily\cyrillicfontsf{CMU Sans Serif}
\newfontfamily\cyrillicfonttt{CMU Typewriter Text}

\usepackage{microtype}


% ====== МАТЕМАТИКА ======
\usepackage{amsmath, amssymb, amsthm, mathtools, bm}


% Русские стандарты математической типографики
\renewcommand{\ge}{\geqslant}
\renewcommand{\le}{\leqslant}
\newcommand{\E}{\mathbb{E}}
\renewcommand{\P}{\mathbb{P}}
\newcommand{\F}{\mathcal{F}}

% ====== ГРАФИКА И TIKZ ======
\usepackage{tikz}
\usetikzlibrary{arrows.meta, positioning, automata, backgrounds, decorations.pathmorphing, calc}

% ====== ОФОРМЛЕНИЕ СТРАНИЦЫ ======
\usepackage[left=2cm, right=2cm, top=2.5cm, bottom=2.5cm]{geometry}
\usepackage{parskip} % Красивые отступы между абзацами
\usepackage{xcolor}


% ====== ГИПЕРССЫЛКИ ======
\usepackage[colorlinks=true, linkcolor=defblue!90!black, urlcolor=defblue!90!black, citecolor=probgreen]{hyperref}

% ====== TCOLORBOX (ОКРУЖЕНИЯ ДЛЯ ТЕОРЕМ И ЗАДАЧ) ======
\usepackage[most]{tcolorbox}
\tcbuselibrary{skins, breakable, theorems}

% Цветовая палитра "МЦНМО/ВШЭ"
\definecolor{defblue}{RGB}{20, 75, 140}
\definecolor{thmred}{RGB}{160, 30, 45}
\definecolor{probgreen}{RGB}{25, 110, 55}
\definecolor{warnpurple}{RGB}{130, 30, 130}
\definecolor{insightteal}{RGB}{0, 100, 100}
\definecolor{solgray}{RGB}{90, 90, 90}
\definecolor{codebg}{RGB}{245, 245, 248}

% Определение
\newtcbtheorem[auto counter, number within=section]{definition}{Определение}{
    enhanced, breakable,
    colback=defblue!3, colframe=defblue,
    fonttitle=\bfseries\sffamily, title style={color=defblue},
    attach boxed title to top left={yshift=-2.5mm, xshift=5mm},
    boxed title style={sharp corners},
    boxrule=1pt, drop fuzzy shadow
}{def}

% Теорема
\newtcbtheorem[auto counter, number within=section]{theorem}{Теорема}{
    enhanced, breakable,
    colback=thmred!3, colframe=thmred,
    fonttitle=\bfseries\sffamily, title style={color=thmred},
    attach boxed title to top left={yshift=-2.5mm, xshift=5mm},
    boxed title style={sharp corners},
    boxrule=1pt, drop fuzzy shadow
}{thm}

% Задача
\newtcbtheorem[auto counter, number within=section]{problem}{Задача}{
    enhanced, breakable,
    colback=probgreen!3, colframe=probgreen,
    fonttitle=\bfseries\sffamily, title style={color=probgreen},
    attach boxed title to top left={yshift=-2.5mm, xshift=5mm},
    boxed title style={sharp corners},
    boxrule=1pt, drop fuzzy shadow
}{prob}

% Важное замечание / Парадокс
\newtcolorbox{paradox}[1][]{
    enhanced, breakable,
    colback=warnpurple!4, colframe=warnpurple!90!black,
    title={\textbf{Внимание: #1}},
    fonttitle=\bfseries\sffamily,
    coltitle=black,
    attach title to upper,
    after title={\par\smallskip},
    boxrule=1.5pt, arc=4pt
}

% Олимпиадный инсайт
\newtcolorbox{insight}[1][]{
    enhanced, breakable,
    colback=insightteal!3, colframe=insightteal, leftrule=4pt,
    boxrule=0pt, arc=0pt,
    title={\textbf{Олимпиадный инсайт: #1}},
    coltitle=insightteal, fonttitle=\bfseries\sffamily,
    attach title to upper,
    after title={\par\smallskip}
}

% Решение
\newenvironment{solution}{
    \begin{tcolorbox}[
        enhanced, breakable, blanker,
        borderline west={2.5pt}{0pt}{solgray},
        left=12pt, top=5pt, bottom=5pt,
        before skip=10pt, after skip=15pt
    ]
    \textbf{\textsf{РЕШЕНИЕ.}}\;
}{
    \hfill $\blacksquare$
    \end{tcolorbox}
}



\begin{document}

\begin{center}
    \vspace*{1.0cm}
    {\Huge \bfseries Матрица Судеб\par}
    \vspace{0.4cm}
    {\Large \itshape Искусство Цепей Маркова без решения СЛАУ\par}
    \vspace{0.6cm}
    \rule{0.6\textwidth}{1pt}\par
    \vspace{0.6cm}
    {\large \textbf{Калашников Александр}}\par
    \vspace{0.2cm}
    {\small Wildberries \& Russ, ФКН ВШЭ}
    \vspace{1.0cm}
\end{center}

\tableofcontents
\newpage

\section{Введение: Проклятие деревьев исходов}

В классических задачах на случайные процессы (генерация текстов LLM, блуждания по графам, игры с разорением) школьники обычно используют \textbf{Метод первых шагов}. Они вводят ожидаемое время $E_i$ для каждого состояния, составляют Систему Линейных Алгебраических Уравнений (СЛАУ) и пытаются решить её вручную.

Если состояний 3, это терпимо. Но если их 5, 10 или 50? Решать СЛАУ методом Гаусса на бумаге — это гарантированная ошибка в арифметике. Более того, возня с дробями скрывает от нас \textit{внутреннюю геометрию} процесса. 

Сегодня мы переведём случайную величину на язык высшей алгебры. Мы узнаем, как предсказывать будущее системы на миллион шагов вперёд одним умножением, и как знаменитая школьная формула бесконечной убывающей геометрической прогрессии $S = \frac{1}{1-q}$ элегантно обобщается на многомерные пространства.

\section{Анатомия Цепи: Матрица Переходов}

\begin{definition}{Цепь Маркова}{markov}
Случайный процесс называется Цепью Маркова, если он обладает \textbf{отсутствием памяти} (Марковским свойством): вероятность перехода в следующее состояние зависит \textit{только от текущего} и не зависит от истории того, как система в него пришла.
\end{definition}

Зададим $n$ состояний. Пусть $p_{ij}$ — вероятность того, что система из состояния $i$ за ровно один шаг перейдёт в состояние $j$. Упакуем все эти вероятности в квадратную матрицу $P$ размера $n \times n$.

\begin{insight}[Свойство Матрицы $P$ (Стохастичность)]
Сумма элементов в \textbf{каждой строке} матрицы $P$ обязана быть строго равна $1$. 
Почему? Строка $i$ описывает все возможные пути из состояния $i$. Так как система обязана куда-то перейти (или хотя бы остаться на месте с вероятностью $p_{ii}$), сумма шансов образует полную группу событий, то есть $1$.
\end{insight}

Обозначим через вектор-строку $\bm{\pi}_0 = (p_1, p_2, \dots, p_n)$ наше начальное распределение (с какими вероятностями мы находимся в состояниях в момент времени $0$).
Распределение системы через 1 шаг находится матричным умножением: $\bm{\pi}_1 = \bm{\pi}_0 P$.

\begin{theorem}{Матрица Будущего}{future}
Распределение вероятностей состояний ровно через $k$ шагов равно:
\[ \bm{\pi}_k = \bm{\pi}_0 P^k \]
В олимпиадном программировании это позволяет находить точные вероятности для $k=10^{18}$ за $\mathcal{O}(\log k)$ с помощью алгоритма бинарного возведения матрицы в степень!
\end{theorem}

\newpage
\section{Бесконечность: Стационарное распределение}

Представьте, что бот бесконечно блуждает по ссылкам в интернете. Куда он придёт через $k \to \infty$ шагов? 
Теорема Перрона — Фробениуса гласит: если граф сильно связен (из любой вершины можно дойти до любой) и апериодичен, система рано или поздно \textbf{забудет своё начальное состояние} $\bm{\pi}_0$ и сойдётся к единому глобальному равновесию.

Это равновесие называется \textbf{Стационарным распределением} $\bm{\pi}$. Раз оно равновесное, то при применении к нему ещё одного шага (умножения на матрицу $P$), оно не должно измениться! Получаем элегантное уравнение:
\[ \bm{\pi} P = \bm{\pi} \]

\begin{paradox}[Собственные векторы и Спектральная Щель]
Посмотрите на уравнение $\bm{\pi} P = 1 \cdot \bm{\pi}$. 
Это в точности определение \textbf{левого собственного вектора} матрицы $P$, соответствующего собственному значению $\lambda_1 = 1$. Линейная алгебра гарантирует, что у любой стохастической матрицы такой вектор всегда существует.

А как быстро система забудет старт? Это зависит от \textbf{второго по модулю} собственного значения $|\lambda_2| < 1$. Величина $\gamma = 1 - |\lambda_2|$ называется \textbf{спектральной щелью} (Spectral Gap). Чем она больше, тем быстрее граф «перемешивает» вероятности. На этом свойстве построены алгоритмы MCMC и экспандерные графы.
\end{paradox}

\begin{problem}{PageRank: Самая дорогая матрица в истории}{}
Интернет состоит из 3 страниц: A, B, C. 
Из A есть ссылки на B и C. Из B ссылка только на C. Из C ссылки на A и B.
Пользователь кликает по случайной ссылке на текущей странице. 
Какую долю времени он будет проводить на странице C в пределе?
\end{problem}
\begin{solution}
Алгоритм Google PageRank — это буквально поиск стационарного вектора $\bm{\pi}$ матрицы интернета! Составим граф и матрицу $P$ (строки A, B, C):

\begin{minipage}{0.4\textwidth}
\begin{center}
\begin{tikzpicture}[thick, >=Stealth, scale=0.9]
  \node[circle, draw=defblue, fill=defblue!10, minimum size=1.0cm, font=\bfseries] (A) at (0, 0) {A};
  \node[circle, draw=thmred, fill=thmred!10, minimum size=1.0cm, font=\bfseries] (B) at (4, 0) {B};
  \node[circle, draw=probgreen, fill=probgreen!10, minimum size=1.0cm, font=\bfseries] (C) at (2, -3.46) {C};

  \path[->] (A) edge[bend left=15] node[above right] {1/2} (B);
  \path[->] (A) edge[bend right=15] node[above left] {1/2} (C);
  \path[->] (B) edge[bend left=15] node[right] {1} (C);
  \path[->] (C) edge[bend left=15] node[left] {1/2} (A);
  \path[->] (C) edge[bend right=15] node[below right] {1/2} (B);
\end{tikzpicture}
\end{center}
\end{minipage}%
\begin{minipage}{0.6\textwidth}
\[ P = \begin{pmatrix} 0 & 1/2 & 1/2 \\ 0 & 0 & 1 \\ 1/2 & 1/2 & 0 \end{pmatrix} \]
Решаем систему $\bm{\pi} P = \bm{\pi}$ (умножая строку на столбцы):
\begin{cases}
    \pi_A = \frac{1}{2}\pi_C \\
    \pi_B = \frac{1}{2}\pi_A + \frac{1}{2}\pi_C \\
    \pi_C = \frac{1}{2}\pi_A + 1\pi_B
\end{cases}
\end{minipage}

Из первого: $\pi_C = 2\pi_A$. Подставляем во второе: $\pi_B = 0.5\pi_A + 0.5(2\pi_A) = 1.5\pi_A$.
Помним условие нормировки вероятностей: $\pi_A + \pi_B + \pi_C = 1$.
Тогда $\pi_A + 1.5\pi_A + 2\pi_A = 1 \implies 4.5\pi_A = 1 \implies \pi_A = 2/9$. \\
Вектор $\bm{\pi} = \left(\mathbf{\frac{2}{9}, \frac{3}{9}, \frac{4}{9}}\right)$. Пользователь проводит $\approx 44\%$ времени на странице C. 
\end{solution}

\newpage
\section{Олимпиадный Чит-код: Неориентированные графы}

Решать СЛАУ $\bm{\pi} P = \bm{\pi}$ для графа на 100 вершин нереально. Но олимпиадники обожают давать задачи на блуждания по \textit{неориентированным} графам (где дороги двусторонние). Здесь кроется гениальный математический шорткат.

\begin{insight}[Принцип Детального Равновесия]
В физике и теории графов существует понятие обратимости. Если граф неориентированный, поток вероятности из вершины $u$ в $v$ в стационаре в точности равен потоку из $v$ в $u$: $\pi_u p_{uv} = \pi_v p_{vu}$. 
Так как шаг делается случайно, $p_{uv} = \frac{1}{\deg(u)}$. Подставляя это, получаем $\frac{\pi_u}{\deg(u)} = \frac{\pi_v}{\deg(v)} = \text{const}$. Это означает, что вероятности прямо пропорциональны степеням вершин!
\end{insight}

\begin{theorem}{Стационарность неориентированных графов}{graph}
В случайном блуждании по связному \textbf{неориентированному} графу доля времени, проведённая в вершине $v$, прямо пропорциональна её степени $\deg(v)$:
\[ \bm{\pi_v = \frac{\deg(v)}{\sum \deg(u)}} = \frac{\deg(v)}{2 |E|} \]
А \textbf{ожидаемое время возвращения} (через сколько шагов бот впервые вернётся в $v$, если стартовал из $v$) равно обратной величине: $\bm{\E[\tau_{v \to v}] = \frac{1}{\pi_v}}$.
\end{theorem}

\begin{problem}{Прогулка Короля (Эргодическая цепь)}{king}
Шахматный король бесконечно долго делает случайные разрешенные ходы по пустой доске $8 \times 8$. 
1. Какую долю времени он проводит в 4 центральных клетках (d4, d5, e4, e5)? \\
2. Если король стоит в углу (a1), через сколько ходов он в среднем вернётся в a1?
\end{problem}
\begin{solution}
Доска — это граф из 64 вершин. Степень — количество возможных ходов.
\begin{itemize}
    \item 4 угла: степень 3. Сумма = $4 \times 3 = 12$.
    \item 24 клетки на краях (без углов): степень 5. Сумма = $24 \times 5 = 120$.
    \item 36 внутренних клеток: степень 8. Сумма = $36 \times 8 = 288$.
\end{itemize}
Общая сумма степеней (удвоенное число рёбер графа) $\sum \deg = 12 + 120 + 288 = \mathbf{420}$.

1) Центральных клеток 4, их суммарная степень равна $4 \times 8 = \mathbf{32}$.
Ответ: Король находится в центре $\mathbf{\frac{32}{420} = \frac{8}{105}} \approx 7.62\%$ времени. 

2) Степень угла a1 равна 3. Стационарная вероятность $\pi_{a1} = \frac{3}{420} = \frac{1}{140}$.
Время возврата равно перевернутой дроби: $\E[\tau_{a1 \to a1}] = \mathbf{140}$ ходов. 
\end{solution}

\begin{paradox}[Ловушка Четности: Конь против Короля]
А сработала бы формула $32/420$ для \textbf{Коня}? И да, и нет. 
Конь при каждом ходе меняет цвет клетки. Граф ходов коня — \textbf{двудольный}. 

Из-за этого строгий математический предел $\lim_{n \to \infty} \bm{\pi}_0 P^n$ \textbf{не существует} (вероятности осциллируют: на четных шагах конь 100\% на белых клетках, на нечетных 100\% на черных). 
Однако, \textit{усредненная доля времени} (предел по Чезаро) и время возврата всё равно идеально подчиняются формуле! Жюри ВсОШ обожает снимать баллы за отсутствие оговорки про двудольность.
\end{paradox}

\newpage
\section{ШЕДЕВР: Поглощающие цепи и Матрица $N$}

Стационарное распределение ищет бесконечный предел. Но в олимпиадах игра чаще всего \textbf{заканчивается} (муха долетела до паука, игрок разорился, сгенерировалась нужная строка). Такие финальные состояния называются \textbf{поглощающими} ($p_{ii} = 1$). 

Для решения таких задач существует феноменальная \textbf{Блочная Теорема}.

\subsection{Разрезаем матрицу}
Запишем матрицу $P$ в каноничном блочном виде. Сначала идут $t$ «транзитных» (временных) состояний, затем $r$ «поглощающих» (финишей).
\[
P = \left( \begin{array}{c|c}
Q & R \\ \hline
0 & I
\end{array} \right)
\]
Где $Q$ — квадратная матрица $t \times t$, описывающая переходы \textit{строго между временными} состояниями, а $R$ — вероятности \textit{прыжка в финиш}.

\subsection{Ряд Неймана (Матричная геометрическая прогрессия)}
Посчитаем ожидаемое количество визитов в каждое временное состояние.
Благодаря линейности матожидания, ожидаемое число визитов — это сумма вероятностей оказаться там на каждом шаге $k=0, 1, 2 \dots$

В момент $0$ мы находимся в старте (единичная матрица $I$). На шаге $1$ нас описывает матрица $Q$. На шаге $2$ — матрица $Q^2$. Полное ожидаемое число посещений — это сумма бесконечного матричного ряда: $N = I + Q + Q^2 + Q^3 + \dots$

Вспомним школьную формулу суммы убывающей геометрической прогрессии: $1 + q + q^2 + \dots = \frac{1}{1-q} = (1-q)^{-1}$. Работает ли это для матриц? \textbf{ДА!}
Доказательство:
\[ (I - Q)(I + Q + Q^2 + \dots) = I + Q + Q^2 \dots - Q - Q^2 - Q^3 \dots = I \]
Умножив на обратную матрицу, получаем $(I - Q)^{-1} = I + Q + Q^2 + \dots$

\begin{theorem}{Фундаментальная Матрица (Матрица Времени)}{fundamental}
Ожидаемое количество посещений каждого состояния до поглощения задаётся \textbf{Фундаментальной матрицей} $N$:
\[ \mathbf{N = (I - Q)^{-1}} \]
Элемент $N_{ij}$ показывает, сколько раз в среднем мы посетим состояние $j$, если начали в $i$.

Чтобы найти \textbf{общее ожидаемое время игры} $E_i$ из состояния $i$, нужно просто сложить все элементы $i$-й строки: $\mathbf{\vec{E} = N \cdot \mathbf{1}}$ (где $\mathbf{1}$ — вектор из единиц).
\end{theorem}

\begin{insight}[Инверсия $(I-Q)$ — это и есть Метод Первых Шагов!]
Если вы запишете систему уравнений времени $\vec{E}$ для графа, вы получите: $\vec{E} = \mathbf{1} + Q \vec{E}$.
Перенесем $\vec{E}$ влево: $\vec{E} - Q \vec{E} = \mathbf{1} \implies (I - Q)\vec{E} = \mathbf{1}$.
Умножив слева на обратную матрицу, получим $\vec{E} = (I - Q)^{-1} \mathbf{1}$. 
Вся "магия" марковских цепей — это просто элегантная матричная запись решения классической СЛАУ!
\end{insight}

\newpage
\section{Парадокс Конвея: Матрицы "читают" тексты}

В методичке по Мартингалам мы обсуждали знаменитый парадокс генерации строк. Монета бросается до выпадения слов $W_1=$ \texttt{HT} (Орёл-Решка) или $W_2=$ \texttt{HH} (Орёл-Орёл). Одинаково ли среднее время ожидания? 

Давайте посмотрим, как Фундаментальная матрица расправляется с этой задачей без всяких искусственных инвариантов. Транзитные состояния автомата — это наибольший префикс искомого слова, который мы уже набрали.

\begin{problem}{Ожидание строки HT}{}
Генерируем \texttt{H} и \texttt{T} с вероятностью $1/2$. Ждем появление \texttt{HT}.
Транзитные состояния: $\emptyset$ (пусто), \texttt{H}. Поглощающее: \texttt{HT}.
\end{problem}
\begin{solution}
Составим матрицу $Q$ (строки/столбцы $\emptyset$, \texttt{H}):
\begin{itemize}
    \item Из $\emptyset$: при \texttt{T} остаемся в $\emptyset$ (1/2), при \texttt{H} идем в \texttt{H} (1/2).
    \item Из \texttt{H}: при \texttt{H} топчемся на месте в \texttt{H} (1/2), при \texttt{T} уходим в поглощение (исчезаем из $Q$).
\end{itemize}
\[ Q = \begin{pmatrix} 1/2 & 1/2 \\ 0 & 1/2 \end{pmatrix} \implies I - Q = \begin{pmatrix} 1/2 & -1/2 \\ 0 & 1/2 \end{pmatrix} \]
Определитель $\det(I-Q) = 1/4$. Обращаем матрицу $2 \times 2$:
\[ N = (I-Q)^{-1} = 4 \begin{pmatrix} 1/2 & 1/2 \\ 0 & 1/2 \end{pmatrix} = \mathbf{\begin{pmatrix} 2 & 2 \\ 0 & 2 \end{pmatrix}} \]
Время из старта $\emptyset$ (сумма 1-й строки): $2 + 2 = \mathbf{4}$ броска.
\end{solution}

\begin{problem}{Ожидание строки HH}{}
Транзитные состояния: $\emptyset$ (пусто), \texttt{H}. Поглощающее: \texttt{HH}.
\end{problem}
\begin{solution}
Матрица $Q$ (строки/столбцы $\emptyset$, \texttt{H}):
\begin{itemize}
    \item Из $\emptyset$: при \texttt{T} остаемся в $\emptyset$ (1/2), при \texttt{H} идем в \texttt{H} (1/2).
    \item Из \texttt{H}: при \texttt{T} \textbf{сбрасываемся в $\emptyset$} (1/2), при \texttt{H} уходим в поглощение.
\end{itemize}
\[ Q = \begin{pmatrix} 1/2 & 1/2 \\ 1/2 & 0 \end{pmatrix} \implies I - Q = \begin{pmatrix} 1/2 & -1/2 \\ -1/2 & 1 \end{pmatrix} \]
Определитель $\det = 1/2 - 1/4 = 1/4$. Обращаем матрицу:
\[ N = (I-Q)^{-1} = 4 \begin{pmatrix} 1 & 1/2 \\ 1/2 & 1/2 \end{pmatrix} = \mathbf{\begin{pmatrix} 4 & 2 \\ 2 & 2 \end{pmatrix}} \]
Время из старта $\emptyset$ (сумма 1-й строки): $4 + 2 = \mathbf{6}$ бросков.
\end{solution}

\begin{insight}[Матрица "видит" структуру слова]
Посмотрите на элемент $N_{11}$ (ожидаемое количество возвратов в $\emptyset$). Для \texttt{HT} оно равно 2, а для \texttt{HH} оно равно 4! Матрица автоматически учла тот факт, что неудачный бросок (\texttt{T}) при ожидании \texttt{HH} отбрасывает нас в самое начало, тогда как при ожидании \texttt{HT} мы можем безопасно "топтаться" на букве \texttt{H} (ведь \texttt{TH} заканчивается на \texttt{H}). 
Фундаментальная матрица $N$ содержит в себе всю теорию детерминированных автоматов и префикс-функций в чистом алгебраическом виде!
\end{insight}


\newpage
\section{Уничтожение задач уровня ВсОШ}

Применим матричное оружие к зубодробительной задаче с регионального этапа. Мы не просто решим её, но и \textbf{вскроем комбинаторику переходов}, которая обычно ставит всех в тупик.

\begin{problem}{Три мухи на пятиугольнике (ВсОШ)}{flies}
Три мухи одновременно садятся на правильный пятиугольник — в 3 различные вершины. Каждую минуту \textit{каждая} муха случайно и независимо перелетает в одну из 4 других вершин (отличных от текущей). 
Найдите ожидаемое время $\E[T]$, когда \textbf{все три мухи впервые окажутся в одной вершине}.
\end{problem}

\begin{solution}
У трёх мух $4^3 = 64$ возможных варианта прыжка. Разобьем состояния на 3 макросостояния:
\begin{itemize}
    \item \textbf{S3:} Все мухи в разных вершинах (старт).
    \item \textbf{S2:} Две мухи в одной вершине, одна отдельно.
    \item \textbf{S1:} Все мухи в одной вершине (\textbf{Поглощающее состояние}).
\end{itemize}

\textbf{Шаг 1. Вскрытие комбинаторики (Из S3).} Пусть мухи сидят в $A, B, C$.
\begin{itemize}
    \item \textbf{В S1:} Они должны прыгнуть в одну общую вершину. Прыгать в $A, B, C$ всем вместе нельзя (кто-то нарушит правило и останется на месте). Значит, они летят либо все в $D$, либо все в $E$. $\implies \mathbf{p_{31} = 2/64}$.
    
    \item \textbf{В S3:} Все 3 цели $(v_1, v_2, v_3)$ должны быть \textbf{различны}. Всего способов выбрать 3 различные цели из 5: $5 \times 4 \times 3 = 60$. Но нам нужно вычесть те, где кто-то остался на месте.
    Применим \textbf{Формулу Включений-Исключений (PIE)}. Свойства $P_1, P_2, P_3$ — муха осталась на месте. 
    Один на месте: если $v_1=A$, две другие мухи выбирают 2 различные цели из оставшихся 4 вершин $\implies 4 \times 3 = 12$. Таких свойств три $\implies \Sigma = 36$.
    Двое на месте: если $v_1=A, v_2=B$, третья муха выбирает из 3 вершин $\implies 3$. Таких пар три $\implies \Sigma = 9$.
    Трое на месте: $v_1=A, v_2=B, v_3=C \implies 1$.
    Число валидных исходов: $60 - 36 + 9 - 1 = \mathbf{32}$. $\implies \mathbf{p_{33} = 32/64}$.
    
    \item \textbf{В S2 (остаток):} $64 - 2 - 32 = 30$ способов. $\implies \mathbf{p_{32} = 30/64}$.
\end{itemize}

\textbf{Шаг 2. Вскрытие комбинаторики (Из S2).} Пусть мухи сидят в $A, A, B$.
\begin{itemize}
    \item \textbf{В S1:} Прыгают все вместе в свободные $C, D$ или $E$ (3 способа). $\implies \mathbf{p_{21} = 3/64}$.
    
    \item \textbf{В S3:} Снова ищем $(v_1, v_2, v_3)$ различные (60 вариантов). Вычитаем нарушения ($v_1=A \lor v_2=A \lor v_3=B$).
    Один на месте: $|P_1| = 12, |P_2| = 12, |P_3| = 12$. Сумма: $36$.
    Двое на месте: $|P_1 \cap P_2| = 0$ (цели обязаны быть различны, они не могут обе прыгнуть в A!). 
    $|P_1 \cap P_3| = 1 \times 3 \times 1 = 3$. $|P_2 \cap P_3| = 3$. Сумма пар: $0 + 3 + 3 = 6$.
    Трое на месте $\implies 0$.
    Число валидных исходов: $60 - 36 + 6 - 0 = \mathbf{30}$. $\implies \mathbf{p_{23} = 30/64}$.
    
    \item \textbf{В S2 (остаток):} $64 - 3 - 30 = 31$ способ. $\implies \mathbf{p_{22} = 31/64}$.
\end{itemize}

\textbf{Шаг 3. Фундаментальная Матрица.}
Выписываем матрицу $Q$ (индексы S3, S2) и вычитаем из $I$:
\[ I - Q = \begin{pmatrix} 1 & 0 \\ 0 & 1 \end{pmatrix} - \frac{1}{64} \begin{pmatrix} 32 & 30 \\ 30 & 31 \end{pmatrix} = \frac{1}{64} \underbrace{\begin{pmatrix} 32 & -30 \\ -30 & 33 \end{pmatrix}}_{M} \]

Для обращения матрицы воспользуемся линейным свойством $(kM)^{-1} = \frac{1}{k} M^{-1}$. Значит $(I-Q)^{-1} = 64 \cdot M^{-1}$.
Формула для $2 \times 2$: $M^{-1} = \frac{1}{\det M} \begin{pmatrix} d & -b \\ -c & a \end{pmatrix}$.
Определитель $\det(M) = 32 \cdot 33 - (-30)(-30) = 1056 - 900 = \mathbf{156}$.

Тогда Фундаментальная матрица $N$:
\[ N = 64 \cdot \frac{1}{156} \begin{pmatrix} 33 & 30 \\ 30 & 32 \end{pmatrix} = \mathbf{\frac{16}{39} \begin{pmatrix} 33 & 30 \\ 30 & 32 \end{pmatrix}} \]

Нам нужно ожидаемое время старта из состояния S3 (первая строка матрицы $N$). Просто суммируем элементы первой строки:
\[ \E[T] = \frac{16}{39} (33 + 30) = \frac{16 \cdot 63}{39} = \frac{16 \cdot 21}{13} = \mathbf{\frac{336}{13}} \]
Без громоздких систем уравнений и подстановок мы получили идеальный ответ $336/13 \approx 25.85$ минут!
\end{solution}

\newpage
\section{Кем мы будем в конце? (Вероятности исходов)}

Иногда нас интересует не \textit{время} игры, а \textit{исход}. Если поглощающих состояний несколько (например, «Игрок победил» или «Разорился»), как узнать вероятность закончить в каждом из них?

\begin{theorem}{Матрица Итога}{absorption}
Вероятность того, что стартовав в транзитном состоянии $i$, система погибнет в конкретном поглощающем состоянии $j$, задаётся матрицей $B$:
\[ \mathbf{B = N \cdot R = (I - Q)^{-1} R} \]
Здесь $N$ описывает, \textit{сколько раз} мы побываем на «трамплинах» перед финалом, а $R$ — вероятности самого финального смертельного прыжка с этих трамплинов в цель.
\end{theorem}

\begin{problem}{Игра с разорением (Связь с Мартингалами)}{ruin}
У игрока 1 рубль. Он выигрывает 1 рубль с вероятностью $p$ и проигрывает с $q$. Игра заканчивается, когда у него 0 (Разорение) или 3 рубля (Победа). Найти вероятность победы матричным методом.
\end{problem}
\begin{solution}
Транзитные состояния: 1, 2. Поглощающие: 0, 3.
Матрица $Q = \begin{pmatrix} 0 & p \\ q & 0 \end{pmatrix} \implies I - Q = \begin{pmatrix} 1 & -p \\ -q & 1 \end{pmatrix}$.
Определитель $\det(I-Q) = 1 - pq$.
Фундаментальная матрица: $N = \frac{1}{1-pq} \begin{pmatrix} 1 & p \\ q & 1 \end{pmatrix}$.

Матрица прыжков в финал $R$ (строки: из 1 и 2; столбцы: в 0 и 3). 
Из 1 в 0 шанс $q$, в 3 шанс $0$. Из 2 в 0 шанс $0$, в 3 шанс $p$.
$R = \begin{pmatrix} q & 0 \\ 0 & p \end{pmatrix}$.

Матрица исходов $B = N \cdot R$:
\[ B = \frac{1}{1-pq} \begin{pmatrix} 1 & p \\ q & 1 \end{pmatrix} \begin{pmatrix} q & 0 \\ 0 & p \end{pmatrix} = \frac{1}{1-pq} \begin{pmatrix} q & \mathbf{p^2} \\ q^2 & p \end{pmatrix} \]
Смотрим первую строку (старт из 1 рубля) и второй столбец (состояние 3 — Победа): вероятность равна $\mathbf{\frac{p^2}{1-pq}}$.

\textbf{Абсолютный триумф алгебры:} В методичке по Мартингалам мы выводили формулу победы $\frac{(q/p) - 1}{(q/p)^3 - 1}$. 
Распишем разность кубов: $\frac{r - 1}{r^3 - 1} = \frac{1}{r^2 + r + 1}$. Подставим $r = q/p$:
\[ \frac{1}{(q/p)^2 + q/p + 1} = \frac{p^2}{q^2 + pq + p^2} \]
Заметим, что так как $p+q=1 \implies (p+q)^2 = 1 \implies p^2+2pq+q^2 = 1 \implies q^2+pq+p^2 = 1-pq$.
Подставляем это в знаменатель и получаем ровно $\mathbf{\frac{p^2}{1-pq}}$. Матрицы и Мартингалы дали абсолютно идентичный результат. Цикл замкнулся!
\end{solution}

\newpage

\begin{tcolorbox}[colback=insightteal!5, colframe=insightteal, title=\textbf{Почему это важно для ИИ? (Reinforcement Learning)}, fonttitle=\bfseries]
В классических Цепях Маркова вероятности $P$ фиксированы законами физики или правилами игры. Но что если мы дадим агенту (например, нейросети) \textbf{право выбора} действий $a$ в каждом состоянии? 

Тогда матрица $P$ превращается в функцию от действий: $P(a)$. Процесс становится \textbf{Марковским Процессом Принятия Решений (MDP)}. Вся область Reinforcement Learning (включая алгоритм RLHF, который сделал ChatGPT умным) — это алгоритмы поиска такой матрицы действий $P(a)$ (называемой Policy), которая максимизирует итоговую награду $V = (I - \gamma P)^{-1} R$, где $\gamma$ — discount factor. Знаменитое Уравнение Беллмана — это в точности наша матрица $N$!
\end{tcolorbox}

\vspace{0.4cm}
\begin{tcolorbox}[colback=codebg, colframe=defblue!80!black, title=\textbf{\texttt{/} Шпаргалка ML-Разработчика: Цепи Маркова в NumPy}, fonttitle=\bfseries\sffamily, boxrule=1pt]
В олимпиадах по ИИ (и в промышленном Data Science) матрицу $(I-Q)^{-1}$ руками не считают. Python решает эту задачу в 4 строки, даже если у вас 1000 состояний.
\vspace{0.2cm}

\fontfamily{qcr}\selectfont
\textcolor{thmred}{import} numpy \textcolor{thmred}{as} np

\textcolor{gray}{\# 1. Поиск ожидаемого времени (Поглощающая цепь)}\\
Q = np.array([[32/64, 30/64], [30/64, 31/64]])\\
\textcolor{gray}{\# Используем linalg.solve (численно стабильнее, чем inv)}\\
times = np.linalg.solve(np.eye(\textcolor{thmred}{len}(Q)) - Q, np.ones(\textcolor{thmred}{len}(Q)))\\
\textcolor{thmred}{print}(times[0]) \textcolor{gray}{\# Ответ: 25.84615...}

\vspace{0.2cm}
\textcolor{gray}{\# 2. Поиск Стационарного распределения (Эргодическая цепь)}\\
P = np.array([[0, 0.5, 0.5], [0, 0, 1], [0.5, 0.5, 0]])\\
eigenvalues, eigenvectors = np.linalg.eig(P.T) \textcolor{gray}{\# Левый вектор = P.T}\\
pi = eigenvectors[:, np.isclose(eigenvalues, 1)]\\
pi = pi / np.sum(pi) \textcolor{gray}{\# Нормировка вектора в сумму 1}\\
\textcolor{thmred}{print}(pi.real.flatten()) \textcolor{gray}{\# Ответ: [0.222, 0.333, 0.444]}
\normalfont
\vspace{0.2cm}

Этот код работает за $\mathcal{O}(V^3)$ и легко вскрывает NLP задачи на генерацию токенов, где классическое ДП убило бы вас сложностью реализации.
\end{tcolorbox}

\end{document}